\begin{proof}
Throughout, Theorem \claimref{synthesis.hadamard-dirac-states-scale-to-the-flat-massless-vacuum-at-every-point-of-a-gravitational-field} is cited as (H): its clause (i) gives the form $U/\sigma + V\log\sigma$ with $U = \Delta^{1/2}R$ and $D_{x}\sigma = \hb^{-1}(\dot\gamma(1))$, the physical tangent at $x$ of the geodesic $\gamma$ from $x'$ to $x$, and its clause (ii) the Taylor data (labelled computed there). The form of the parametrix is Fact \claimref{synthesis.imported-hadamard-parametrix-of-the-squared-dirac-operator-leads-with-van-vleck-root-times-spin-transport}. The curvature operator is $\mathcal{R}(B) = -R(\hb(B))$, $R(a\wedge b) = a\cdot\nabla\Omega(b) - b\cdot\nabla\Omega(a) + \Omega(a)\times\Omega(b)$ the covariant field strength (Definition \claimref{obj.synthesis.gtg-flat-spacetime-gauge-fields}(c)); (H)(ii) is written with $\mathcal{R}$.

\emph{(i).} $D_{x}(U/\sigma) = (D_{x}\sigma^{-1})U + \sigma^{-1}D_{x}U = -sU/\sigma^{2} + (D_{x}U)/\sigma$, so the leading part of $G(x, x')$ is $-\Delta^{1/2}sR/\sigma^{2}$. Exchanging the points reverses the geodesic: $\sigma$ and $\Delta$ are symmetric, the rotor becomes its reverse, and $D_{x'}\sigma$ is the physical tangent at $x'$ of the reversed geodesic. The physical tangent is transported by the rotor, $\hb^{-1}(\dot\gamma(1)) = R\,\hb^{-1}(\dot\gamma(0))\,\rev{R}$ (the geodesic equation and $\dot R = -\tfrac12\Omega(\dot\gamma)R$), so $s' = -\rev{R}sR$ and $Q(s') = Q(s) = 2\sigma$. Then
\[
\langle a\,(sR)\,b\,(s'\rev{R})\rangle_{0} = -\langle a\,s\,R\,b\,\rev{R}\,s\,R\,\rev{R}\rangle_{0} = -\langle a\,s\,\hat b\,s\rangle_{0} = -\bigl(2(a\cdot s)(\hat b\cdot s) - (a\cdot\hat b)Q(s)\bigr),
\]
by $R\rev{R} = 1$ and the trace $\langle a\,x\,b\,x\rangle_{0} = 2(a\cdot x)(b\cdot x) - (a\cdot b)Q(x)$ of clause (i) of Theorem \claimref{synthesis.one-loop-vacuum-polarisation-extends-without-a-regulator-and-ward-preserving-extensions-differ-only-by-the-renormalisation-constants}. Multiplying by $\Delta^{1/2}\Delta^{1/2}/\sigma^{4}$ gives $\Delta\,\Pi^{0}$. In $\langle a\,(sR)\,b\,\rev{R}\rangle_{0}$ and $\langle a\,R\,b\,(s'\rev{R})\rangle_{0}$ the argument is an odd element, with no scalar part; and $\langle a\,R\,b\,\rev{R}\rangle_{0} = a\cdot\hat b$. The parametrix of (H) carries the field theory's unit on one of the two terms $D_{x}$ and $m$; stripping it leaves the parities and the traces unchanged but not the relative sign of the mass-mass and leading terms, which involves the square of the unit. The check verifies the identities with four exact rotors (Cayley transforms of rational bivectors, simple and not), six vectors $s$ and sixteen pairs $a$, $b$, as a control, and that the identity fails if the tangent at $x'$ is not transported.

\emph{(ii).} The statement $a\cdot D_{x}R = \tfrac14\mathcal{R}(a\wedge\xi)R + O(\xi^{2})$ is covariant under rotation gauge transformations, so it may be verified in one gauge. Take $x_{0}$ with $\hb(x_{0}) = 1$, first derivatives of $\hb$ zero at $x_{0}$, and
\[
\Omega(a)(x_{0} + z) = -\tfrac12\,\mathcal{R}(z\wedge a) + O(z^{2}) .
\]
Its field strength at $x_{0}$ is $a\cdot\nabla\Omega(b) - b\cdot\nabla\Omega(a) = -\tfrac12\mathcal{R}(a\wedge b) + \tfrac12\mathcal{R}(b\wedge a) = -\mathcal{R}(a\wedge b)$, which is the stored relation between the field strength and the curvature operator; any $\Omega$ vanishing at $x_{0}$ with this field strength differs from it at first order by the symmetric part of $a\cdot\nabla\Omega(b)$, which a rotation gauge transformation quadratic in $z$ removes. By (H)(ii), with $x = x_{0} + ly$, $x' = x_{0} + ly'$, $d = y - y'$, $\bar y = \tfrac12(y + y')$: the transport rotor is $1 - \tfrac{l^{2}}{2}(\bar y\cdot\nabla)\Omega(d) + O(l^{3}) = 1 + \tfrac{l^{2}}{4}\mathcal{R}(\bar y\wedge d) + O(l^{3})$. The scaled derivative $l\,a\cdot D = a\cdot\partial_{y} + \tfrac{l}{2}\Omega(a)(x_{0} + ly) + O(l^{2})$ applied to it gives, at order $l^{2}$,
\[
\tfrac14\bigl[\tfrac12\mathcal{R}(a\wedge d) + \mathcal{R}(\bar y\wedge a)\bigr] - \tfrac14\mathcal{R}(y\wedge a) = \tfrac14\bigl[\tfrac12\mathcal{R}(a\wedge d) - \tfrac12\mathcal{R}(d\wedge a)\bigr] = \tfrac14\mathcal{R}(a\wedge d),
\]
using $y - \bar y = \tfrac12 d$; with $\xi = ld$ this is the claim, and it does not depend on $\bar y$. Contracting, $D_{x}R = \tfrac14\gamma^{\mu}\mathcal{R}(e_{\mu}\wedge\xi)R$, and $\gamma^{\mu}\mathcal{R}(e_{\mu}\wedge\xi) = \partial_{a}\cdot\mathcal{R}(a\wedge\xi) + \partial_{a}\wedge\mathcal{R}(a\wedge\xi) = \mathcal{R}(\xi)$ by the protraction identity, Theorem \claimref{synthesis.protraction-identity-of-the-curvature-operator-from-vanishing-torsion}(iii), the field being torsion-free. A curvature operator of Weyl type is by definition one with $\partial_{a}\mathcal{R}(a\wedge\xi) = 0$ for every $\xi$. By (H)(ii), $\Delta = 1 + \tfrac16\Ric^{\circ}(\xi, \xi) + O(\xi^{3})$ with $\Ric^{\circ}(\xi, \xi) = -\xi\cdot\mathcal{R}(\xi)$, so $\Delta^{1/2} = 1 - \tfrac1{12}\xi\cdot\mathcal{R}(\xi)$ and its gradient is $-\tfrac16\mathcal{R}(\xi)$ (the Ricci map is symmetric, by the same clause). Hence $D_{x}U = (-\tfrac16 + \tfrac14)\mathcal{R}(\xi)R + O(\xi^{2}) = \tfrac1{12}\mathcal{R}(\xi)R + O(\xi^{2})$. For $\mathcal{R}(B) = \tfrac{\Lambda}{3}B$, $\mathcal{R}(\xi) = \Lambda\xi$ and $D_{x}U = \tfrac{\Lambda}{12}\xi R$. The corrections from $\hb - 1 = O(z^{2})$ and from the deviation of the geodesic from the straight line are of relative order $l^{2}$ and enter $D_{x}R$ at order $\xi^{2}$. The check verifies that the gauge has field strength $-\mathcal{R}$ and that the gauge of the opposite sign does not, the derivative formula for a general linear $\mathcal{R}$ at exact points (the symmetric difference of a quadratic polynomial is its derivative), the contraction on operators of type D ($B \mapsto c(B + 3\sigma_{r}B\sigma_{r})$ with $\sigma_{r}^{2} = 1$), on sums of rotated copies and on operators $a\wedge b \mapsto \tfrac12(P(a)\wedge b + a\wedge P(b))$ of Ricci type, that a map violating the protraction identity leaves a trivector, and the de Sitter value, against which the opposite relative sign gives $-\tfrac{5}{12}\Lambda\xi$.

\emph{(iii) (argued).} Write $G(x, x') = -\Delta^{1/2}sR/\sigma^{2} + (D_{x}U)/\sigma + mU/\sigma + Vs/\sigma + (D_{x}V + mV)\log\sigma$. With the rules (SD1) to (SD4) of Fact \claimref{synthesis.imported-epstein-glaser-extension-theorem-and-ward-preserving-time-ordered-products}(vi), in the normal coordinates of $x'$: the first term has degree $3$; $mU/\sigma$ degree $2$; $(D_{x}U)/\sigma$ degree at most $1$ by (ii), with coefficient the Ricci vector map; $Vs/\sigma$ degree at most $1$, with coefficient the coincidence value of $V$; the logarithmic term degree at most $0$. In the product: degree $6$ is $\Pi^{0}$; degree $5$ could only come from the leading part times the mass term, which vanishes by (i); degree $4$ comes from mass times mass, from the leading part times $(D_{x}U)/\sigma$, from the leading part times $Vs/\sigma$, and from $(\Delta - 1)\Pi^{0}$ with $\Delta - 1 = \tfrac16\Ric^{\circ}(\xi, \xi) + O(\xi^{3})$; the coincidence value of $V$ is by the first transport equation a multiple of the squared Dirac operator applied to $U$ at coincidence, where the second derivative of the rotor enters through the contraction of (ii) and the zero-order term is a scalar by the Lichnerowicz form. The Weyl part is absent from all of these coefficients. It enters at degree $3$ through first derivatives of the curvature operator in $D_{x}U$ and $\Delta$, and at degree $2$ through the term $\tfrac1{180}\operatorname{tr}K_{\xi}^{2}$ of $\Delta$ (cited), for which the check verifies on the Weyl-type samples that the tidal map has zero trace and nonzero trace of its square. This clause is argued because the transport equation is not stored and the rules are imported from memory; the higher-order Hadamard form needed for the remainder is the rule (SD4), the stored parametrix theorem giving only class $C^{1}$.

\emph{(iv) (argued).} By (i), for fixed $x'$, in the normal coordinates $\xi'$ of $x'$ ($s' = D_{x'}\sigma$ is minus the position vector of $x$ in them) and with $a$ transported to $x'$, $\Pi^{0} = -(2(\hat a\cdot s')(b\cdot s') - (\hat a\cdot b)Q(s'))/\sigma^{4}$ is the flat massless kernel, homogeneous of degree $-6$ in $\xi'$. Given (iii), Fact \claimref{synthesis.imported-epstein-glaser-extension-theorem-and-ward-preserving-time-ordered-products}(ii) applied to the time-ordered kernel in $\xi'$ gives $\omega = 6 - 4 = 2$ and that two extensions of degree $6$ differ by $\sum_{|\alpha| \le 2}c_{\alpha}\partial^{\alpha}\delta$; there are $1 + 4 + 10 = 15$ multi-indices. The coefficients $c_{\alpha}$ are chosen for each $x'$ separately; as functions of $x'$ they are unconstrained by this argument. For any smooth $c$, $\partial_{\mu}(c\,F^{\mu\nu})$ has zero divergence, so the transverse local terms on a fixed field include $c(x)$ times the flat transverse term for every $c$. The stored flat extension of clause (ii) of the cited vacuum-polarisation theorem, read in $\xi'$, extends $\Pi^{0}$ with the stored logarithmic coefficient. By Fact \claimref{synthesis.imported-local-covariance-with-scaling-of-time-ordered-products}(iii), if the extension is taken local and covariant with almost homogeneous scaling, then in the difference of two such extensions $c_{\alpha}$ with $|\alpha| = j$ is a polynomial with constant coefficients in the curvature operator, its covariant derivatives and the mass, of dimension $2 - j$: a constant at $j = 2$, nothing at $j = 1$ beyond terms that are total derivatives of those, and at $j = 0$ a combination of $m^{2}$ and the curvature operator. This removes the dependence on $x'$ as a free function. Theorem \claimref{synthesis.curved-field-vacuum-polarisation-freedom-is-five-local-covariant-constants-of-which-one-is-transverse}, which depends on this one and is cited here for what follows from it, not as a premise, counts the list (five constants) and shows that two such extensions that are both transverse differ by $c\,D\cdot(D\wedge A)$, the single constant $Z_{3}$. The existence of a transverse local covariant extension needs the Ward identity on the field, which is not derived here and is not part of that Fact. The first Hadamard transport equation is used in (iii) for one coefficient at degree $4$ only, and not in this clause.
\end{proof}
